\section{Cost Amortization and Distortions}
\label{sec:cost_results}
Rollup L1 publication costs exhibit dual-layer characteristics: a fixed cost associated with submitting the batch and verifying the validity proof, and a variable cost that scales with transaction data size. Increasing batch sizes allows the sequencer to amortize the fixed costs over more transactions, reducing the per-transaction L1 gas fee.
In the Stage 1 sweeps (utilizing a Mock Verifier contract with a low ~29,000 gas overhead), the true weighted gas per transaction scaled as follows:
\begin{itemize}
    \item \textbf{Batch Size 25 (s1\_bs\_0025):} True Weighted Gas/Tx was \textbf{20,549 gas}.
    \item \textbf{Batch Size 100 (s1\_bs\_0100):} True Weighted Gas/Tx was \textbf{19,681 gas}.
    \item \textbf{Batch Size 500 (s1\_bs\_0500):} True Weighted Gas/Tx was \textbf{19,459 gas}.
\end{itemize}
\subsection{Empty Batch Cost Distortion}
A significant statistical anomaly was identified in the unweighted aggregation framework. For the large batch run \texttt{s1\_bs\_1000}, the reported gas cost was \textbf{83,626 gas/tx}, whereas smaller batches reported ~19,600 gas. 
Telemetry logs showed that the run produced 6 batches: Batches 1 and 2 were full (605 and 661 txs), Batch 3 had 1 tx, and Batches 4, 5, and 6 had 0 txs (empty). Because the metrics collector uses a simple mean of the ratios and applies a floor of $\max(\text{tx\_count}, 1)$ to prevent division by zero, the empty batches were calculated as:
\begin{equation}
\text{Gas/Tx}_{\text{empty}} = \frac{112,332\text{ gas}}{1\text{ tx}} = 112,332\text{ gas/tx}
\end{equation}
Averaging these ratios yields the distorted 83.6k gas/tx. The true weighted average (total gas divided by total transactions) was actually \textbf{19,771 gas/tx}, matching the amortization curve. This underscores the necessity of using true weighted averages when reporting rollup transaction economics.
